\section*{Three-node workload comparison: calibration and protocol}
The three-node Fig.~5 experiment uses NJ-k1 as cloud and NJ-k2/NJ-k3 as request origins with local execution capacity. All three nodes are KVM virtual machines running Debian~13, each exposing 16 logical CPUs. NJ-k1/NJ-k2 use Xeon Platinum 8163 processors; NJ-k3 uses a Xeon Gold 6442Y. The allocated application CPUs are 0--3 on NJ-k1 and 0--1 on each end-side node. There are two/one/one prestarted serial execution slots, each with a two-CPU quota, 512-MiB memory limit and Go parallelism set to one. New experiment control processes and clients use CPUs 8--15. The matched CPU and memory limits apply to the Go application workers; framework control processes, including Ray's Python actors and routers, run outside those worker limits. Unrelated host services remain present. All four experiment deployments remain resident, with workload traffic applied to one system at a time.

\begin{table}[htbp]
\centering\small
\caption{Pre-comparison calibration medians in milliseconds. Each entry uses 20 sequential samples after three unscored warmups per kernel. Bare columns report kernel service time; Tide columns report the complete local request path.}
\label{tab:three-node-calibration}
\begin{tabular}{lrrrrrr}
\toprule
& \multicolumn{2}{c}{CPU} & \multicolumn{2}{c}{Memory-access} & \multicolumn{2}{c}{Synchronous I/O}\\
Node & Bare & Tide & Bare & Tide & Bare & Tide\\
\midrule
NJ-k1 & 21.66 & 32.53 & 82.67 & 92.37 & 57.20 & 67.31\\
NJ-k2 & 20.93 & 31.40 & 82.23 & 91.09 & 61.91 & 71.53\\
NJ-k3 & 4.17 & 9.44 & 43.88 & 50.15 & 11.41 & 17.25\\
\bottomrule
\end{tabular}
\end{table}

For workload $w$, the fixed standalone reference is $C_w=2/s_{1w}+1/s_{2w}+1/s_{3w}$, where $s_{iw}$ is the median Bare kernel service time in seconds. The original six calibration files are retained unchanged: CPU $379.9193$, memory-access $59.1457$ and synchronous I/O $138.7640$ requests/s. The complete new Fig.~5 measures nominal fractions 10\%,20\%,\ldots,100\%, with the same integer rule $\operatorname{round}(\max(2,pC_w/100))$. In increasing load order, CPU rates are 38/\allowbreak76/\allowbreak114/\allowbreak152/\allowbreak190/\allowbreak228/\allowbreak266/\allowbreak304/\allowbreak342/\allowbreak380, memory rates are 6/\allowbreak12/\allowbreak18/\allowbreak24/\allowbreak30/\allowbreak35/\allowbreak41/\allowbreak47/\allowbreak53/\allowbreak59, and I/O rates are 14/\allowbreak28/\allowbreak42/\allowbreak56/\allowbreak69/\allowbreak83/\allowbreak97/\allowbreak111/\allowbreak125/\allowbreak139 requests/s. These are total rates across both origins. The capacity scale is fixed before comparative outcomes and does not predict framework saturation throughput. Fresh diagnostic calibrations are preserved separately and do not retune rates or execution estimates. The 30\%/80\% points are newly measured in the same randomized sweep, not imported from the earlier campaign. In particular, slot-count round robin does not compensate for the faster NJ-k3 processor.

Tide's node-specific execution estimates are the local-path medians rounded up to whole milliseconds: CPU 33/32/10, memory-access 93/92/51, and I/O 68/72/18 for NJ-k1/NJ-k2/NJ-k3. All nodes use a 500-ms deadline, equal score weights $(1,1,1)$, 100-ms latency scale, 5000-ms maximum state age, 1-ms configured per-node network estimate and 24-hour Stub TTL. The TTL is an idle-time limit, refreshed by accepted/completed work, rather than a maximum process lifetime; the cloud role does not start the idle-expiry timer. The sixteen preformal worker identities remain unchanged across this campaign.

All four frameworks use the same Go executable and workload inputs. The CPU workload applies 100 SHA-256 iterations to 64~KiB. The memory-access workload performs four sequential read/modify/write passes over 16~MiB; it is not assumed to be purely memory-bound. Each I/O request creates its own temporary file and performs four 4-MiB write, sync, readback and CRC32-check cycles, then removes the file. Bare, Tide, Ray and KubeEdge use the same host ext4 directory on each host VM; mount device/inode checks verify this identity. Checksums and processed-byte counts are independently computed and checked for every returned result.

The Ray baseline uses Ray~2.58.0 with Python~3.13.5, four fixed Serve replicas, two logical CPUs per replica and one ongoing request per replica. Native Serve routing retains local-node preference and may send work to other nodes. The KubeEdge baseline uses real EdgeCore~1.23.1 on the two end-side hosts and cloud-side Kubernetes workers. Its dedicated control plane has no shared overlay routing; application Pods use host networking and explicit experiment ingress. Clients cycle the four KubeEdge/Bare slots, with a two-slot phase offset at NJ-k3. Thus the baselines share worker inventory and slot-count routing, while Tide and Ray use their native placement paths. Each Ray replica forwards requests over local HTTP to its Go worker. The KubeEdge cloud ingress adds a dedicated TCP relay to the nested cloud node, and edge ingress uses container port publication. These transport adapters are included in the measurements; the comparison concerns the integrated end-to-end paths.

The common driver opens one HTTP connection per offer, without retries, a client concurrency cap or a serial client queue. Both origins use a common future wall-clock start; preformal records show all three hosts reporting NTP synchronization, without bounding cross-host clock error or proving continuous synchronization; measured per-origin latency uses local monotonic clocks. The 500-ms deadline is measured from actual client submission, while planned-arrival-to-return time and replayer lateness are retained separately, together with clock-mapping drift. Burst synchronization is nominal and subject to host clock and client scheduling error. No queueing time is inferred by subtracting timestamps from different hosts.

Each 10-s cell uses balanced, 90\%-NJ-k2 hotspot or one-second burst arrivals. Five paired repetitions randomize all 90 workload/rate/pattern cells and the four systems within each cell (seed 20261025), yielding 1,800 treatments and 3,600 origin cohorts. A treatment collects responses before the next one. A timeout or uncertain server-side failure stops further execution until a recorded quiescence check; such outcomes remain in the data. The check verifies unchanged worker identities and observes worker CPU, I/O, threads and relevant sockets after the response-timeout horizon; Ray queue and in-flight gauges are also retained. This is bounded observed quiescence. The frozen native interface does not expose a live admission-queue counter, so the check does not establish an exact count of pending native routines. The raw files retain all offers, failures, successful latency, destination, worker identity and kernel service time. Run-level Student-$t$ intervals use five repetitions and four degrees of freedom, without clipping interval bounds.

The native harness preserves actual container execution and cross-node messaging. Validation identified and corrected an unknown-Stub synchronization update that could overwrite the first known Stub's state. The external measurement ingress also registers the originating request and invokes the native receipt/result-wait protocol for remote execution; cloud handlers are registered for the prestarted worker inventory. These corrections were tested before the frozen comparative campaign. The older single-host measurements use a different deployment and collection protocol and remain separately labeled in Tables~S3--S4; numerical differences between those runs and Fig.~5 cannot be attributed to topology alone.

During the earlier 30\%/80\% study, an initial campaign stopped after fourteen treatments when its cloud TCP relay reached a 1024-descriptor soft limit (37 connection resets). Those diagnostic records remain separate. The new ten-level campaign starts with the established HTTP/control/collector limit of at least 65,536 descriptors and a separately frozen randomized order. Native PAIX workers retain their original 1024-descriptor limit (seven open descriptors in preformal evidence); their message-queue execution path does not use one HTTP connection per offer. This distinction does not change worker CPU or memory allocations.

\paragraph{Complete ten-level campaign outcomes and timing}
The independent audit verifies all 1,800 treatments, 3,600 origin cohorts and 1,907,400 offers, with 476,850 offers per system. All 269,015 failed offers are recorded timeouts: Bare 23,046, KubeEdge 21,678, Tide 82,213 and Ray 142,078. No failed request is removed from the offered denominator or replayed. All sixteen original workers remain represented. The complete sweep tables report all 360 goodput cells, each with five run-level repetitions and pointwise, non-simultaneous $t$(4) intervals. Earlier Table~\ref{tab:three-node-goodput} retains the independently dated 30\%/80\% campaign and is not the source of the new Fig.~5 curves.

Among checksum-valid successful returns, the cloud/NJ-k2/NJ-k3 destination fractions are 31.85/17.71/50.44\% for Tide, 32.70/20.82/46.48\% for Ray, 48.26/25.49/26.25\% for Bare and 48.22/25.60/26.17\% for KubeEdge. Timeouts have unknown destinations. These fractions describe returned work, not routing shares of all offers, and do not isolate a scheduling mechanism.

The largest run-level replayer-lateness P95 and P99 are 128.658 and 219.972~ms; the largest individual lateness is 319.810~ms. These are maxima across runs, not pooled percentiles. The maximum local clock-mapping drift is 0.048~ms; it is not a cross-host clock-offset bound. Scheduling delays can affect arrival shape and fixed-window goodput. Successful-latency statistics exclude timeouts. Fixed-window goodput counts checksum-valid returns inside each origin's half-open 10-s window that also meet the inclusive 500-ms actual-submit-to-return deadline. Successful returns after that window remain in the cohort records but not the plotted goodput. Five repetitions on one three-VM allocation with these kernels and pinned framework versions limit generalization.

The frozen campaign is \texttt{f6-1792867033070427000}; all six measurement-source hashes, reference calibrations, traces, runtime identities, request records and failure/drain evidence are retained. A post-measurement DNS failure prevented installation of remote plotting packages after the remote independent audit had passed. Local raw-data collection and a second independent audit precede local plotting; this recovery does not launch or replace any treatment.

\paragraph{Earlier single-host mixed-deadline protocol}
Supplementary Table~S4 retains the earlier NJ-k1 CPU experiment, independently of the three-node protocol above. Native Tide SubmitRoutine/PAIX, a direct Go HTTP worker, one serial Ray Serve~2.58.0 replica and one KubeEdge~1.23.1 EdgeCore-managed Pod in isolated kind use the same calibrated Go CPU kernel and two-logical-CPU worker budget. Ray permits one ongoing request; the KubeEdge Pod uses fixed \texttt{nodeName} placement with requests/limits of two CPUs and 512~MiB. Placement and replica counts remain fixed, startup is excluded, and memory limits and control-process budgets follow the respective deployments. Tide uses a 22-ms CPU execution estimate, 2000-ms Stub TTL, equal score weights $(1,1,1)$, 100-ms latency scale and 5000-ms maximum node-state age. These estimates and TTL do not apply to the three-node experiment.

Each total offered rate of 20/40/80 requests/s has three 6-s runs, alternating 100/500-ms deadlines so that each class receives half the offers. On-time goodput divides successful timely returns within the window by 6~s; yield denominators include all submissions and successful latency includes response drain. Table~S4 retains every three-run class-specific goodput mean and its run-level 95\% confidence interval. The original single-host figure assets remain archived; Supplementary Fig.~S6 reports the new three-node experiment with its own calibration and protocol. The earlier experiment characterizes the combined single-host admission, queueing and framework path. Changes in deployment and collection protocol prevent attributing numerical differences from Fig.~5 solely to topology.
